% Class-compatibility preamble retained for a reader excerpt of
%   Sahar Diskin, Anna Geisler, "Perfect Matching in Product Graphs and in their
%   Random Subgraphs", arXiv:2404.14020v4, CC BY 4.0.
%
% The source pins its own publisher class (scrartcl); the delivery wrapper accepts
% only the classes article and amsart, so this file supplies the amsart baseline and
% then copies the author's own declaration block mechanically and unchanged
% (pinned-original-source.tex, source lines 30-57: the author's newtheorem
% declarations and the author's macro declarations).  The author's own package
% list of source lines 4-21 is reused only for the entries the retained text needs
% (enumitem, amsmath, amssymb, amsthm, mathtools, xcolor, geometry) and the
% cross-reference packages already needed by the author's own \Cref/\eqref usage
% (hyperref, cleveref); tikz, caption, subcaption, float, graphicx, wrapfig and
% abstract serve figures and the title block that this excerpt does not carry.
% No mathematical macro, statement, symbol or assumption is introduced, changed
% or removed by this file.
\documentclass[10pt,a4paper]{amsart}
\usepackage[shortlabels]{enumitem}
\usepackage{amsmath}
\usepackage{amssymb}
\usepackage{amsthm}
\usepackage{mathtools}
\usepackage{xcolor}
\usepackage[margin=1in]{geometry}
\usepackage{hyperref}
\usepackage{cleveref}
\hypersetup{
  colorlinks=true,
  linkcolor=blue,
  filecolor=magenta,
  urlcolor=cyan
}

\newtheorem{theorem}{Theorem}
\newtheorem{thm}{Theorem}[section]
\newtheorem{corollary}[thm]{Corollary}
\newtheorem{lemma}[thm]{Lemma}
\newtheorem{proposition}[thm]{Proposition}
\newtheorem{claim}[thm]{Claim}
\newtheorem{remark}[thm]{Remark}
\newtheorem{question}[thm]{Question}

\newcommand{\whp}{whp}
\newcommand{\defi}{\operatorname{def}}
\newcommand{\odd}{\operatorname{odd}}
\newcommand{\iso}{\operatorname{iso}}
\newcommand{\cN}{\mathcal{N}}
\newcommand{\cO}{\mathcal{O}}
\newcommand{\cE}{\mathcal{E}}
\newcommand{\cK}{\mathcal{K}}
\newcommand{\cP}{\mathcal{P}}
\newcommand{\cQ}{\mathcal{Q}}

\newcommand{\footremember}[2]{%
  \footnote{#2}%
  \newcounter{#1}%
  \setcounter{#1}{\value{footnote}}%
}

\newcommand{\ag}[1]{\textcolor{purple}{(Anna: #1)}}
\newcommand{\sd}[1]{\textcolor{red}{(Sahar: #1)}}
\crefname{originallemma}{original lemma}{original lemmas}
\Crefname{originallemma}{Original lemma}{Original lemmas}
\makeatletter
\newlabel{lem:ge-boundary-identity}{{7.4}{source}{}{sourcepointer-1}{}}
\newlabel{lem:ge-boundary-identity@cref}{{[originallemma][0][]7.4}{[1][0][]source}{}{}{}}
\makeatother
\crefname{originallemma}{original lemma}{original lemmas}
\Crefname{originallemma}{Original lemma}{Original lemmas}
\makeatletter
\newlabel{lem:strict-hall}{{7.2}{source}{}{sourcepointer-2}{}}
\newlabel{lem:strict-hall@cref}{{[originallemma][0][]7.2}{[1][0][]source}{}{}{}}
\makeatother
\crefname{originallemma}{original lemma}{original lemmas}
\Crefname{originallemma}{Original lemma}{Original lemmas}
\makeatletter
\newlabel{lem:two-colour-count}{{7.7}{source}{}{sourcepointer-3}{}}
\newlabel{lem:two-colour-count@cref}{{[originallemma][0][]7.7}{[1][0][]source}{}{}{}}
\makeatother
\crefname{originalequation}{original equation}{original equations}
\Crefname{originalequation}{Original equation}{Original equations}
\makeatletter
\newlabel{eq:colouring-weight}{{7.3}{source}{}{sourcepointer-4}{}}
\newlabel{eq:colouring-weight@cref}{{[originalequation][0][]7.3}{[1][0][]source}{}{}{}}
\makeatother
\crefname{originalequation}{original equation}{original equations}
\Crefname{originalequation}{Original equation}{Original equations}
\makeatletter
\newlabel{eq:ge-exact}{{7.4}{source}{}{sourcepointer-5}{}}
\newlabel{eq:ge-exact@cref}{{[originalequation][0][]7.4}{[1][0][]source}{}{}{}}
\makeatother
\crefname{originalequation}{original equation}{original equations}
\Crefname{originalequation}{Original equation}{Original equations}
\makeatletter
\newlabel{eq:strong-3d}{{7.5}{source}{}{sourcepointer-6}{}}
\newlabel{eq:strong-3d@cref}{{[originalequation][0][]7.5}{[1][0][]source}{}{}{}}
\makeatother
\crefname{originalequation}{original equation}{original equations}
\Crefname{originalequation}{Original equation}{Original equations}
\makeatletter
\newlabel{eq:expansion-log}{{1.1}{source}{}{sourcepointer-7}{}}
\newlabel{eq:expansion-log@cref}{{[originalequation][0][]1.1}{[1][0][]source}{}{}{}}
\makeatother
\crefname{originallemma}{original lemma}{original lemmas}
\Crefname{originallemma}{Original lemma}{Original lemmas}
\makeatletter
\newlabel{lem:product-connectivity}{{2.2}{source}{}{sourcepointer-8}{}}
\newlabel{lem:product-connectivity@cref}{{[originallemma][0][]2.2}{[1][0][]source}{}{}{}}
\makeatother
\crefname{originallemma}{original lemma}{original lemmas}
\Crefname{originallemma}{Original lemma}{Original lemmas}
\makeatletter
\newlabel{lem:rooted-trees}{{2.3}{source}{}{sourcepointer-9}{}}
\newlabel{lem:rooted-trees@cref}{{[originallemma][0][]2.3}{[1][0][]source}{}{}{}}
\makeatother
\crefname{originallemma}{original lemma}{original lemmas}
\Crefname{originallemma}{Original lemma}{Original lemmas}
\makeatletter
\newlabel{lem:strong-boundary}{{6.7}{source}{}{sourcepointer-10}{}}
\newlabel{lem:strong-boundary@cref}{{[originallemma][0][]6.7}{[1][0][]source}{}{}{}}
\makeatother
\crefname{originallemma}{original lemma}{original lemmas}
\Crefname{originallemma}{Original lemma}{Original lemmas}
\makeatletter
\newlabel{lem:small-boundary-sets}{{6.3}{source}{}{sourcepointer-11}{}}
\newlabel{lem:small-boundary-sets@cref}{{[originallemma][0][]6.3}{[1][0][]source}{}{}{}}
\makeatother
\crefname{originallemma}{original lemma}{original lemmas}
\Crefname{originallemma}{Original lemma}{Original lemmas}
\makeatletter
\newlabel{lem:three_edges}{{6.2}{source}{}{sourcepointer-12}{}}
\newlabel{lem:three_edges@cref}{{[originallemma][0][]6.2}{[1][0][]source}{}{}{}}
\makeatother
\crefname{originalequation}{original equation}{original equations}
\Crefname{originalequation}{Original equation}{Original equations}
\makeatletter
\newlabel{eq:tutte-violation}{{6.1}{source}{}{sourcepointer-13}{}}
\newlabel{eq:tutte-violation@cref}{{[originalequation][0][]6.1}{[1][0][]source}{}{}{}}
\makeatother
\crefname{originaltheorem}{original theorem}{original theorems}
\Crefname{originaltheorem}{Original theorem}{Original theorems}
\makeatletter
\newlabel{th:expansion}{{2.1}{source}{}{sourcepointer-14}{}}
\newlabel{th:expansion@cref}{{[originaltheorem][0][]2.1}{[1][0][]source}{}{}{}}
\makeatother
\crefname{originalequation}{original equation}{original equations}
\Crefname{originalequation}{Original equation}{Original equations}
\makeatletter
\newlabel{eq:size-X}{{7.2}{source}{}{sourcepointer-15}{}}
\newlabel{eq:size-X@cref}{{[originalequation][0][]7.2}{[1][0][]source}{}{}{}}
\makeatother
\title{In the non-bipartite case the set A of the Gallai-Edmonds decomposition stays below the threshold throughout the window}
\author{Sahar Diskin, Anna Geisler}
\date{arxiv:2404.14020v4; CC BY 4.0}
\begin{document}
\maketitle
\noindent\textit{Editorial note.} All delivered mathematical text is the author's own, copied from the pinned author source sources/combinatorics-fl04-vspace-1-5cm-perfect-matching-in-p/pinned-original-source.tex; no statement or proof step was rewritten, reordered or removed.\par\medskip
\noindent\textit{Editorial note.} Preamble provenance: the source pins the publisher class scrartcl, which the delivery wrapper does not accept, so this card uses sources/combinatorics-fl04-vspace-1-5cm-perfect-matching-in-p/retained-preamble.tex, which supplies the amsart baseline and then carries the author's own newtheorem and macro declarations (pinned source lines 30-57) verbatim. No macro was invented and no mathematical symbol was changed.\par\medskip
\noindent\textit{Editorial note.} Bibliography provenance: the source carries no thebibliography environment and ends with bibliographystyle abbrv plus bibliography perc (author source lines 1691-1692), so the bibliography is resolved with BibTeX from the identical copy staged with the same arXiv version under sources/product-graph-matchings-2404.14020v4/perc.bib; no entry was edited and the bibliography style is the author's own abbrv.\par\medskip
\noindent\textit{Editorial note.} The counted proof is the longest in the paper and relies on the author's lemmas on the exact boundary identity (source lines 1099-1103), on the strict Hall inequality (source lines 1073-1079) and on the weighted counting bound for bipartitions (source lines 1344-1364); all three have their own cards in this lane, so here they are navigation pointers only and their content is declared as named standard prerequisites.\par\medskip
\noindent\textit{Editorial note.} The threshold u\_0 that the statement compares against is delivered verbatim as a context span, quoted from the author's own display.\par\medskip
\noindent\textit{Allowed standard prerequisites (editorial).}\begin{itemize}
\item The Gallai-Edmonds decomposition (D, A, C), its parameters m, a, r and h, and the auxiliary bipartite graph Q between A and the components of the D part. Set up in the author's subsection on the decomposition (source lines 1054-1098). Fixes every object the proof manipulates during the counting.
\item The exact edge-counting identity for the Gallai-Edmonds parameters. The author's lemma at source lines 1099-1103, whose proof is delivered by this lane's separate card for that statement. Turns the assumed size of A into an identity relating the boundary of D, the internal edges of A and the component sizes.
\item The strict Hall inequality for the auxiliary bipartite graph, and the weighted counting bound for bipartitions with a small boundary. The author's lemmas at source lines 1073-1079 and 1344-1364, both delivered by this lane's separate cards. Control respectively the matching structure around A and the number of bipartitions of the relevant set, which together give the probabilistic contradiction.
\end{itemize}\medskip
\ifdefined\hypertarget\hypertarget{sourcepointer-1}{}\fi
\noindent\textit{Source reference.} Original lemma 7.4: sources/combinatorics-fl04-vspace-1-5cm-perfect-matching-in-p/pinned-original-source.tex, source lines 1099--1099. Navigation pointer only: the counted proof uses this exact boundary identity in its main estimate; the identity's own statement and proof are delivered by this lane's separate card.\par\medskip
\ifdefined\hypertarget\hypertarget{sourcepointer-2}{}\fi
\noindent\textit{Source reference.} Original lemma 7.2: sources/combinatorics-fl04-vspace-1-5cm-perfect-matching-in-p/pinned-original-source.tex, source lines 1073--1073. Navigation pointer only: the counted proof uses this strict Hall inequality for the auxiliary bipartite graph; its own card exists in this lane and its text is not reproduced here.\par\medskip
\ifdefined\hypertarget\hypertarget{sourcepointer-3}{}\fi
\noindent\textit{Source reference.} Original lemma 7.7: sources/combinatorics-fl04-vspace-1-5cm-perfect-matching-in-p/pinned-original-source.tex, source lines 1344--1344. Navigation pointer only: the counted proof uses this weighted counting bound for bipartitions in its final step; its own card exists in this lane.\par\medskip
\ifdefined\hypertarget\hypertarget{sourcepointer-4}{}\fi
\noindent\textit{Source reference.} Original equation 7.3: sources/combinatorics-fl04-vspace-1-5cm-perfect-matching-in-p/pinned-original-source.tex, source lines 1360--1360. Navigation pointer only: it labels the displayed weighted consequence of the counting lemma that the counted proof applies at the end; the display belongs to that other statement and is not reproduced here.\par\medskip
\ifdefined\hypertarget\hypertarget{sourcepointer-5}{}\fi
\noindent\textit{Source reference.} Original equation 7.4: sources/combinatorics-fl04-vspace-1-5cm-perfect-matching-in-p/pinned-original-source.tex, source lines 1100--1100. Navigation pointer only: the counted proof consumes the estimate or statement labelled here; its content is declared as a named standard prerequisite and the labelled text is not reproduced in this card.\par\medskip
\ifdefined\hypertarget\hypertarget{sourcepointer-6}{}\fi
\noindent\textit{Source reference.} Original equation 7.5: sources/combinatorics-fl04-vspace-1-5cm-perfect-matching-in-p/pinned-original-source.tex, source lines 814--814. Navigation pointer only: the counted proof consumes the estimate or statement labelled here; its content is declared as a named standard prerequisite and the labelled text is not reproduced in this card.\par\medskip
\ifdefined\hypertarget\hypertarget{sourcepointer-7}{}\fi
\noindent\textit{Source reference.} Original equation 1.1: sources/combinatorics-fl04-vspace-1-5cm-perfect-matching-in-p/pinned-original-source.tex, source lines 224--224. Navigation pointer only: the counted proof consumes the estimate or statement labelled here; its content is declared as a named standard prerequisite and the labelled text is not reproduced in this card.\par\medskip
\ifdefined\hypertarget\hypertarget{sourcepointer-8}{}\fi
\noindent\textit{Source reference.} Original lemma 2.2: sources/combinatorics-fl04-vspace-1-5cm-perfect-matching-in-p/pinned-original-source.tex, source lines 318--318. Navigation pointer only: the counted proof consumes the estimate or statement labelled here; its content is declared as a named standard prerequisite and the labelled text is not reproduced in this card.\par\medskip
\ifdefined\hypertarget\hypertarget{sourcepointer-9}{}\fi
\noindent\textit{Source reference.} Original lemma 2.3: sources/combinatorics-fl04-vspace-1-5cm-perfect-matching-in-p/pinned-original-source.tex, source lines 231--231. Navigation pointer only: the counted proof consumes the estimate or statement labelled here; its content is declared as a named standard prerequisite and the labelled text is not reproduced in this card.\par\medskip
\ifdefined\hypertarget\hypertarget{sourcepointer-10}{}\fi
\noindent\textit{Source reference.} Original lemma 6.7: sources/combinatorics-fl04-vspace-1-5cm-perfect-matching-in-p/pinned-original-source.tex, source lines 811--811. Navigation pointer only: the counted proof consumes the statement or estimate labelled here; its content is declared as a named standard prerequisite and the labelled text is not reproduced in this card.\par\medskip
\ifdefined\hypertarget\hypertarget{sourcepointer-11}{}\fi
\noindent\textit{Source reference.} Original lemma 6.3: sources/combinatorics-fl04-vspace-1-5cm-perfect-matching-in-p/pinned-original-source.tex, source lines 714--714. Navigation pointer only: the counted proof consumes the statement or estimate labelled here; its content is declared as a named standard prerequisite and the labelled text is not reproduced in this card.\par\medskip
\ifdefined\hypertarget\hypertarget{sourcepointer-12}{}\fi
\noindent\textit{Source reference.} Original lemma 6.2: sources/combinatorics-fl04-vspace-1-5cm-perfect-matching-in-p/pinned-original-source.tex, source lines 703--703. Navigation pointer only: the counted proof consumes the statement or estimate labelled here; its content is declared as a named standard prerequisite and the labelled text is not reproduced in this card.\par\medskip
\ifdefined\hypertarget\hypertarget{sourcepointer-13}{}\fi
\noindent\textit{Source reference.} Original equation 6.1: sources/combinatorics-fl04-vspace-1-5cm-perfect-matching-in-p/pinned-original-source.tex, source lines 660--660. Navigation pointer only: the counted proof consumes the statement or estimate labelled here; its content is declared as a named standard prerequisite and the labelled text is not reproduced in this card.\par\medskip
\ifdefined\hypertarget\hypertarget{sourcepointer-14}{}\fi
\noindent\textit{Source reference.} Original theorem 2.1: sources/combinatorics-fl04-vspace-1-5cm-perfect-matching-in-p/pinned-original-source.tex, source lines 220--220. Navigation pointer only: the counted text refers to the statement or display labelled here; its content is declared as a named standard prerequisite and the labelled text is not reproduced in this card.\par\medskip
\ifdefined\hypertarget\hypertarget{sourcepointer-15}{}\fi
\noindent\textit{Source reference.} Original equation 7.2: sources/combinatorics-fl04-vspace-1-5cm-perfect-matching-in-p/pinned-original-source.tex, source lines 1378--1378. Navigation pointer only: the counted text refers to the statement or display labelled here; its content is declared as a named standard prerequisite and the labelled text is not reproduced in this card.\par\medskip
\begin{lemma}\label{lem:nonbipartite-ge}
Suppose that at least one base graph of $G$ is non-bipartite.
Whp for every $p\in[p_-,p_+]$, if $\defi(G_p)\ge2$, then the set $A$ in the Gallai--Edmonds decomposition of $G_p$ has size at most $u_0$.
\end{lemma}

\begin{proof}
Fix $p\in[p_-,p_+]$, let $(D,A,C)$ be the Gallai--Edmonds decomposition of $G_p$, and suppose that $\defi(G_p)\ge2$ and $|A|>u_0$.
Let $D_0$ be the union of the singleton components of $G_p[D]$, let $D^+$ be the family of non-singleton components of $G_p[D]$, and put
\[
 D_1=\bigcup_{D_i \in D^+} D_i.
\]
Write $a=|A|$, $r=|D^+|+|D_0|-a\ge2$, and recall
\[
 h=\sum_{i<j}e_G(D_i,D_j)+e_G(D,C).
\]
All the edges counted in $h$ are not present in $G_{p_-}$.
Hence, we shall sum $(1-p_-)^h$ over the choices of $A$, $C$, the singleton components of $D$, and the non-singleton components of $D$.

Every $D_i\in D^+$ is a non-singleton factor-critical component, so $|D_i|\ge3$.
By \eqref{eq:ge-exact},
\begin{equation}\label{eq:mixed-exact}
 h=\frac{dr}{2}+e_G(A)+\frac{e_G(C, C^c)}{2}
   +\frac{1}{2}\sum_{D_i\in D^+}\bigl(e_G(D_i, D_i^c)-d\bigr).
\end{equation}
Also 
\begin{align} \label{eq:h-lower-bound}
    h\ge e_G(D_0),
\end{align}
since every edge of $G[D_0]$ joins two distinct components of $G_p[D]$.
Taking three quarters of \eqref{eq:mixed-exact} and one quarter of \eqref{eq:h-lower-bound} gives
\begin{align} \label{eq:combined-LB-h}
     h\ge\frac{3dr}{8}+\frac34e_G(A)+\frac14e_G(D_0)  +\frac38 e_G(C, C^c)
 +\frac38\sum_{D_i\in D^+}\bigl(e_G(D_i, D_i^c)-d\bigr).
\end{align}

Put $k_i=\min\{|D_i|,|D_i^c|\}$ so that $3\le k_i\le n/2$.
If $k_i\le2u_1$, then \eqref{eq:strong-3d} gives $e_G(D_i, D_i^c)\ge2d$.
If $k_i>2u_1$, then $\ln(n/k_i)\ge\ln 2$, and \eqref{eq:expansion-log} gives
\[
 e_G(D_i, D_i^c)\ge\frac{e\ln 2}{M}k_i>2d
\]
since $u_1\gg d$. Thus,
\begin{align} \label{eq:apply-D_i}
    \frac18\bigl(e_G(D_i, D_i^c)-d\bigr)\ge\frac1{16}e_G(D_i, D_i^c).
\end{align}
Furthermore, 
 \begin{align} \label{eq:replace-Y}
     e_G(C, C^c)+\sum_{D_i\in D^+}e_G(D_i, D_i^c) \geq e_G((C \cup D_1), (C \cup D_1)^c).
 \end{align}
Plugging \eqref{eq:apply-D_i} and \eqref{eq:replace-Y} into \eqref{eq:combined-LB-h} we obtain the following lower bound for $h$:
\begin{equation}\label{eq:mixed-master}
 h\ge\frac{3dr}{8}
 +\frac{e_G(A)+e_G(D_0)}4
 +\frac{e_G((C \cup D_1), (C, \cup D_1)^c)}{16}
 +\frac14\sum_{D_i \in D^+}\bigl(e_G(D_i, D_i^c)-d\bigr).
\end{equation}

We first sum over the possible families $D^+$ for a fixed set $(C \cup D_1)$.
By \Cref{lem:rooted-trees} here are at most $|C \cup D_1|(ed)^{k-1}$ connected sets of cardinality $k$ with a root in $C \cup D_1$.

For $3\le |D_i|\le2u_1$, \Cref{lem:strong-boundary} gives
\begin{equation*}\label{eq:component-small-boundary}
 e_G(D_i, D_i^c)-d \ge d+15(M-1)|D_i|\ln d.
\end{equation*}
Plugging this into the last term of \eqref{eq:mixed-master} the total contribution from these sets is at most
\begin{align*}
    &|C \cup D_1| \sum_{k=3}^{2u_1} (ed)^{k-1} (1-p_-)^{\frac{d+15(M-1)k \ln d}{4}}\\
    \leq & n^{3/4} \omega^{1/4} \sum_{k \geq 3} (ed)^{k-1} d^{-\frac{15 \ln 2}{4} k}\\
    \leq & n^{3/4+o(1)},
\end{align*}
where we used that $\ln n \geq \frac{\ln 2}{M-1} d$ and that $\frac{15 \ln 2}{4}>1$.

For $2u_1<|D_i|\le n/2$, first suppose that
\[
 e_G(D_i, D_i^c) \ge \frac{5(M-1)}{\ln 2}|D_i|\ln d.
\]
Using again $\ln n \geq \frac{\ln 2}{M-1}d$, the contribution from this range is at most
\begin{align*}
    &n (1-p_-)^{-d/4}\sum_{k>2u_1}(ed)^k (1-p_-)^{\frac{5(M-1)}{4 \ln 2} k \ln d}\\
    \leq & \left(\frac{\omega}{n} \right)^{-1/4} \sum_{k >2u_1} (ed)^k d^{-\frac{5}{4}k}
 \le\exp\{-\Omega(u_1\ln d)\}.
\end{align*}
Indeed, $\frac{5}{4}>1$ makes the series geometric with exponent $-\Omega(k\ln d)$, while the prefactor has logarithm $O(d)$ and $u_1\gg d$.

Otherwise if $e_G(D_i, D_i^c)<|D_i|\ln^2 d$, then \Cref{lem:small-boundary-sets} (which is applicable because $2u_1 \geq \frac{n}{d^{\ln^2 d}}$) gives at most $\exp\{2|D_i|/\ln d\}$ choices, whereas
\[
 e_G(D_i, D_i^c)-d = \Omega(|D_i|)
\]
since $|D_i|>2u_1 \gg d$. Thus,
\begin{align*}
    &\sum_{k > 2u_1} \exp\left\{ \frac{2k}{\ln d}\right\} (1-p_-)^{\Omega(k)}\\
    \leq & \sum_{k > 2u_1} \exp\{-\Omega(k)\} = \exp\{-\Omega(u_1)\}.
\end{align*}

There is at most one $D_i$ with $|D_i|>n/2$. We now do a similar argument focusing on its complement $D_i^c$. Since $A\subseteq K^c$, we have $u_0<|D_i^c|<n/2$.
If
\[
 e_G(D_i, D_i^c)\ge \frac{8(M-1)}{\ln(2e)}|D_i^c|\ln(en/|D_i^c|),
\]
then,
\[
 \sum_{u_0<k<n/2}\binom nk (1-p_-)^{\left\{\frac{8(M-1)}{\ln(2e)} k\ln(en/k)-d\right\}/4} =\exp(-\Omega(u_0)),
\]
using that $\log n \geq \frac{\log 2}{M-1} d$ and $2 \log 2>1$.
Otherwise if $e_G(D_i, D_i^c)\le |D_i^c|\ln^2d$, then by \Cref{lem:small-boundary-sets} (which is applicable because $u_0 \geq \frac{n}{d^{\ln^2 d}}$) and $e_G(D_i, D_i^c)-d\ge \Omega(|D_i^c|)$. Here again $|D_i^c|>u_0\gg d$ permits the subtraction of $d$ from the linear expansion bound. Then,
\[
\sum_{u_0<k<n/2}\exp\{2k/\ln d\}(1-p_-)^{\Omega(k)}=\exp(-\Omega(u_0)).
\]

Let $\mathcal{S}_{\le n/2}$ be the sum of the terms corresponding to connected sets of size at most $n/2$, and let $\mathcal{S}_{>n/2}$ be the corresponding sum for a possible connected set larger than $n/2$.
The preceding estimates give
\[
 \mathcal{S}_{\le n/2}\le n^{3/4+o(1)},
 \qquad
 \mathcal{S}_{>n/2}\le e^{-\Omega_M(u_0)}.
\]
A pairwise disjoint family contains at most one member larger than $n/2$.
% Discarding disjointness among the remaining members and using $\prod_K(1+x_K)\le\exp\{\sum_Kx_K\}$ for $x_K\ge0$ gives
% \[
%  (1+\mathcal{S}_{>n/2})\exp\{\mathcal{S}_{\le n/2}\}
%  \le\exp\{n^{3/4+o(1)}\}.
% \]
In particular, uniformly for sets $Y$ with $|V\setminus Y|>u_0$, the sum over unordered families of pairwise disjoint connected subsets of $G[Y]$ satisfies
\begin{equation}\label{eq:component-family-sum}
 \sum_{\substack{D^+\\|V\setminus (C \cup D_1)|>u_0}}
 (1-p_-)^{\frac14\sum_{D_i\in D^+}\{e_G(D_i, D_i^c)-d\}}
 \le\exp\{n^{3/4+o(1)}\}
 =\exp\{o(u_0)\}.
\end{equation}
This sum counts only the non-singleton components of $G_p[D]$; it does not use the vertices or internal edges of $C$. Given $C \cup D_1$, an ordered partition $V\setminus (C \cup D_1)=A\sqcup D_0$, and the family $D^+$, the remaining set
$C=(C \cup D_1)\setminus\bigcup_{D_i\in D^+}D_i$ is determined.

Recalling \eqref{eq:mixed-master}, after \eqref{eq:component-family-sum} and after discarding $(1-p_-)^{3dr/8}\le1$, it remains to bound the following sum over ordered partitions $V=(C \cup D_1)\sqcup A\sqcup D_0$:
\begin{equation}\label{eq:ordered-partition-sum}
  \sum_{\substack{V=(C \cup D_1)\sqcup A\sqcup D_0\\
 |A|>u_0}}  (1-p_-)^{e_G((C \cup D_1), (C \cup D_1)^c)/16+\{e_G(A)+e_G(D_0)\}/4}.  
 %\le e^{-\Omega_M(u_0)}+e^{-\Omega_M(n)}.
\end{equation}
We split according to whether $s \coloneqq |(C \cup D_1)^c|\le n/2$. Suppose first that this inequality holds. Then $s\ge |A|>u_0$ and $\ln\left(\frac{en}{s}\right)=O(\ln d)$.
If
\begin{equation}\label{eq:small-side-high-boundary}
 e_G((C \cup D_1), (C \cup D_1)^c) \ge \frac{20(M-1)}{\ln 2}s\ln\left(\frac{en}{s}\right),
\end{equation}
then
\[
 \binom ns2^s (1-p_-)^{e_G((C \cup D_1), (C \cup D_1)^c)/16}
 \le \exp\{-\Omega(s\ln(en/s))\}.
\]
If \eqref{eq:small-side-high-boundary} fails, then
\[
 e_G((C \cup D_1), (C \cup D_1)^c)=O(s\ln d)<s\ln^2 d.
\]
By \Cref{lem:small-boundary-sets}, there are $e^{o(s)}$ choices for $(C \cup D_1)^c$, while expansion (\Cref{th:expansion}) gives $e_G((C \cup D_1), (C \cup D_1)^c)= \Omega(s)$.
For each fixed $(C \cup D_1)^c$, \Cref{lem:two-colour-count} gives
\[
 \sum_{A\sqcup D_0}  (1-p_-)^{\{e_G(A)+e_G(D_0)\}/4}\le e^{o(s)},
\]
obtained from \eqref{eq:colouring-weight} with $\theta=1/4$.

Thus, the range $s \leq n/2$ contributes at most
\begin{equation}\label{eq:small-side-sum}
 \sum_{s>u_0}e^{-\Omega(s)}=e^{-\Omega(u_0)}.
\end{equation}
In particular, when the Gallai--Edmonds set $C \cup D_1$ contains more than half of vertices, we count the smaller set $A\cup D_0$ rather than $C \cup D_1$.

We may now assume $s=|(C \cup D_1)^c|>n/2$, so $|C \cup D_1|<n/2$.
Fix $\varepsilon>0$ to be chosen sufficiently small. Suppose first that $|C \cup D_1|\ge\varepsilon n$.
If $e_G((C \cup D_1), (C \cup D_1)^c)\ge \frac{20(M-1) \ln 3}{\ln 2}n$, then
\[
 3^n (1-p_-)^{e_G((C \cup D_1), (C \cup D_1)^c)/16}\le e^{-\Omega(n)}.
\]
Otherwise,
\[
 e_G((C \cup D_1), (C \cup D_1)^c)=O(n)=O(|C \cup D_1|)< |C \cup D_1|\ln^2 d.
\]
By \Cref{lem:small-boundary-sets}, there are $e^{o(n)}$ choices for $(C \cup D_1)^c$, and, for each fixed $(C \cup D_1)^c$, \eqref{eq:colouring-weight} with $\theta=1/4$ (which is applicable since $e_G((C \cup D_1), (C \cup D_1)^c)<\frac{20(M-1) \ln 3}{\ln 2}n\le\frac{40(M-1) \ln 3}{\ln 2}|(C \cup D_1)^c|$) gives
\[
 \sum_{R=A\sqcup D_0}
 (1-p_-)^{\{e_G(A)+e_G(D_0)\}/4}\le e^{o(n)}.
\]
Moreover,
\[
 e_G((C \cup D_1), (C \cup D_1)^c)=\Omega(|C \cup D_1|)\ge \alpha_M\varepsilon n,
\]
for some suitable constant $\alpha_M$.
So, in total this contributes $e^{-\Omega(n)}$.

Finally suppose that $|C \cup D_1|<\varepsilon n$.
Fix $i$ such that $H_i$ is non-bipartite.
Fixing all coordinates other than the $i$-th partitions $V$ into $n/|V(H_i)|$ sets, each of which induces a copy of $H_i$.
At most $|C \cup D_1|$ of these copies meet $C \cup D_1$.
In every remaining copy, the vertices lie in $A\cup D_0$; since $H_i$ is non-bipartite, at least one edge has both endpoints in $A$ or both endpoints in $D_0$.
Thus
\begin{equation}\label{eq:nonbipartite-copies}
 e_G(A)+e_G(D_0)\ge n/M-|C \cup D_1|\ge n/(2M).
\end{equation}
The ordered partitions for which $e_G(A)+e_G(D_0)\ge \frac{5(M-1)\log 3}{\log 2}n$ contribute at most
\[
 3^n (1-p_-)^{\frac{5(M-1)\log 3}{4 \log 2} n}=e^{-\Omega(n)}.
\]
For the remaining ordered partitions with $e_G(A)+e_G(D_0)< \frac{5(M-1)\log 3}{\log 2}n$, repeat the counting argument from \Cref{lem:two-colour-count} for the ordered partition $R=A\sqcup D_0$.
The trivial bound gives $e_G((C \cup D_1), (C \cup D_1)^c)\le d|C \cup D_1|$, so \eqref{eq:size-X} gives only $\varepsilon n+O(n/d)$ exceptional vertices in $X$.
As in the proof of \Cref{lem:two-colour-count}, the random sample $T$, together with $A\mathbin{\triangle}N_{G[R]}(T)$ after possibly interchanging $A$ and $D_0$, determines the ordered partition of $R$. We get
\[
|A \triangle N_{G[R]}(T)| \leq \epsilon n + o(n), \qquad |T|=o(n).
\]
The number of choices for $(T, A \triangle N_{G[R]}(T))$ is at most
\[
 \exp\{\varepsilon\ln\left(\frac{e}{\varepsilon} \right)n+o(n)\}
\]
and the choices of $(C \cup D_1)$ contribute at most $\exp\{\varepsilon\ln(e/\varepsilon)n\}$.
Recalling that by \eqref{eq:nonbipartite-copies} there are at least $(e_G(A)+e_G(D_0))/4 \geq n/(8M)$ edges that are not present in $G_{p_-}$. Choosing $\varepsilon=\varepsilon(M)>0$ sufficiently small ensures $2 \varepsilon \log\left( \frac{e}{\varepsilon}\right) < \frac{\log 2}{8M(M-1)}$. Then, the contribution is at most
\begin{align*}
    &\exp\{\varepsilon\ln\left(\frac{e}{\varepsilon} \right)n+o(n)\} (1-p_-)^{n/(8M)}\\
 \leq &\exp\{2 \varepsilon\ln\left(\frac{e}{\varepsilon} \right)n - \frac{\log 2}{8M(M-1)} +o(n)\}\\
 = & \exp(-\Omega(n)).
\end{align*}
Combining \eqref{eq:small-side-sum} with this contribution from $s>\frac{n}{2}$ proves that \eqref{eq:ordered-partition-sum} is at most $\exp(-\Omega(u_0))+ \exp(-\Omega(n))$.

Multiplying \eqref{eq:component-family-sum} with this bound on \eqref{eq:ordered-partition-sum} gives $o(1)$, showing that whp there is no Gallai--Edmonds decomposition of $G_p$ with $\defi(G_p)\ge2$ and $|A|>u_0$.
\end{proof}

\bibliographystyle{plain}
\bibliography{perc}
\end{document}
